% Generated by scripts/txt2tex.py from chapters/chapter_06.txt. Edit the TXT, not this file.
\chapterlayout{06}
\chapterquestion{האם משפט שנאמר בביטחון מלא מגלה משהו על כך שהוא נכון?}
\chapter{מה AI יודע — ומה לא}\label{ch:06}
\section{משפט אחד שאין דרך לבדוק}
התשובה כבר עומדת על המסך, והקוראת עוברת עליה לאט. רוב השורות נשמעות סבירות. ואז היא נעצרת מול משפט אחד, מנוסח בלי היסוס ובלי "ייתכן", שאין לה שום דרך לדעת אם הוא נכון. הטון שלו זהה לטון של כל השאר. שום סימן בשוליים לא מפריד בין מה שהמערכת "יודעת" לבין מה שהיא פשוט ניסחה היטב.
ומחשבה נוספת מתגנבת: היא הקלידה שאלה מהמטבח שלה, בשעה מסוימת, במילים שלה. לאן הלכו המילים האלה? שתי השאלות, האמינות של מה שיוצא מהמערכת וגורלו של מה שנכנס אליה, הן הנושא של הפרק.
\section{כשהטקסט מתנתק מהמקור}
לחוקרים יש שם לתופעה: הזיה (hallucination). לפי סקירה מקיפה של התחום, בגרסתה המעודכנת ב-arXiv, הזיה בעיבוד שפה היא תוכן שנוצר והוא חסר היגיון או אינו נאמן למקור שניתן למערכת \cite{S-ch06-02}. המונח נלקח מהפסיכולוגיה, שם הזיה היא תפיסה לא אמיתית שמרגישה אמיתית \cite{S-ch06-02}. לפי תקציר הסקירה, מערכות שמייצרות טקסט על בסיס למידה עמוקה נוטות להפיק טקסט כזה, והנטייה פוגעת גם בביצועים וגם באמון המשתמשים \cite{S-ch06-01}.
הסקירה מבחינה בין שני סוגים \cite{S-ch06-02}. הזיה פנימית (intrinsic) סותרת את המקור; לדוגמה, סיכום שמייחס למאמר מסקנה הפוכה לזו שבו. הזיה חיצונית (extrinsic) היא תוכן שאי אפשר לאמת מתוך המקור: הוא לא מאושר וגם לא מופרך. בשיחה של הקוראת לא ניתן למערכת שום מקור, ולכן ההגדרה חלה עליה רק בהשאלה. ובכל זאת, בלשון הסקירה, המשפט שעצר אותה קרוב לסוג השני: אולי הוא נכון, אבל דבר בתשובה עצמה אינו מאפשר לאשר אותו או להפריך אותו.
ומאין זה מגיע? לפי הסקירה, הגורמים מתחלקים לשתי משפחות \cite{S-ch06-02}. האחת היא הנתונים: פער בין המקור לבין טקסט הייחוס שעליו המערכת מתאמנת, שנוצר באיסוף נתונים לפי כללי אצבע, או מפני שהמשימה עצמה מזמינה אותו. השנייה היא האימון וההסקה: ייצוג פנימי לא מושלם, פענוח (decoding) שגוי בשלב שבו נבחרות המילים, והטיה של הידע שנאגר בפרמטרים של המודל. בפרק 4 ראינו שהתשובה נבנית בחירה אחר בחירה, וכל בחירה היא גם הזדמנות לסטות.
\section{תיק אחד בניו יורק}
ב-22 ביוני 2023 פרסם בית המשפט הפדרלי למחוז הדרומי של ניו יורק החלטה בתיק Mata v. Avianca \cite{S-ch06-13}. לפי ההחלטה, עורכי דין הגישו לבית המשפט החלטות שיפוטיות שאינן קיימות, עם ציטוטים ואסמכתאות מזויפים שנוצרו בכלי ChatGPT, והמשיכו לעמוד מאחוריהן גם אחרי שבית המשפט העמיד את קיומן בספק \cite{S-ch06-13}. בין היתר מנה בית המשפט בשמן שש החלטות כאלה \cite{S-ch06-13}.
ציר הזמן מלמד יותר מכל הגדרה. ההגשה הראשונה שציטטה את התקדימים הבדויים נעשתה ב-1 במרץ 2023; ב-15 במרץ הטיל הצד שכנגד ספק בקיומם; ב-11 וב-12 באפריל הוציא בית המשפט צווים; והאמת החלה להתגלות רק ב-25 במאי \cite{S-ch06-13}. בסוף הוטל קנס של 5,000 דולר, ובית המשפט חייב את עורכי הדין לשלוח מכתבים לתובע ולכל שופט שזוהה בטעות כמחבר של אחת ההחלטות המדומות \cite{S-ch06-13}.
עם זאת, השופט לא פסל את הטכנולוגיה. בפתח ההחלטה הוא קובע שאין פסול מובנה בשימוש בכלי בינה מלאכותית אמין לשם עזרה, אבל הכללים מטילים על עורכי הדין תפקיד של שומרי סף, שאחראים לדיוק המסמכים שהם מגישים \cite{S-ch06-13}. הכשל לא היה בעצם השימוש, אלא בוויתור על הבדיקה.
\begin{mythbox}
\mythlabel אם מערכת עונה בשטף, אפשר פשוט לשאול אותה אם היא צודקת.
\tcblower
\realitylabel באותו תיק צורפו צילומי מסך של שיחה עם ChatGPT. המשתמש שאל אם התקדימים שקיבל מזויפים, והתשובה הייתה שהם אמיתיים ושאפשר למצוא אותם במאגרים המשפטיים LexisNexis ו-Westlaw \cite{S-ch06-13}. הם לא היו אמיתיים. שאלת הבדיקה קיבלה עוד טקסט שוטף, וטקסט שוטף אינו אימות. את הבדיקה עושים מול מקור חיצוני, לא מול המערכת עצמה.
\end{mythbox}
\section{ביטחון אינו ראיה}
אילו ידעה המערכת לומר "כאן אני לא בטוחה", חלק מהבעיה היה נפתר. לתכונה הזאת קוראים כיול (calibration). במאמר מ-2017 הגדירו גואו ועמיתיו מודל מכויל כמודל שמידת הביטחון שלו תואמת את שיעור ההצלחה שלו בפועל: מתוך 100 תחזיות שבכל אחת מהן הביטחון הוא 0.8, מצופה ש-80 יהיו נכונות \cite{S-ch06-05}.
הממצא שלהם הטריד: רשתות נוירונים מודרניות, בניגוד לרשתות מלפני עשור, מכוילות גרוע \cite{S-ch06-04}. הם השוו רשת בת 5 שכבות מ-1998 (LeNet) לרשת בת 110 שכבות מ-2016 (ResNet) במשימת סיווג תמונות. אצל הוותיקה הביטחון הממוצע היה קרוב לדיוק; אצל החדשה, המדויקת יותר, הוא היה גבוה בהרבה מהדיוק \cite{S-ch06-05}. החוקרים מציינים שהעמקה והרחבה של הרשת עשויות להקטין את השגיאה ובו בזמן לפגוע בכיול, ומדגישים שאינם יכולים לטעון לסיבתיות \cite{S-ch06-05}. הם גם בחנו תיקון חד-פרמטרי, כיול בעזרת טמפרטורה (temperature scaling), שהיה יעיל ברוב מאגרי הנתונים שנבדקו \cite{S-ch06-04}. למרות השם, זו אינה הטמפרטורה מפרק 4, שקובעת כמה אקראית תהיה בחירת המילה הבאה.
אבל המחקר עסק ברשתות שמסווגות תמונות ומסמכים, לא בצ'טבוטים שכותבים פסקאות, ואין בידינו מקור מאומת שמראה שהממצא חל על מודלי שפה גדולים. לכן זו אנלוגיה בלבד, אבל אנלוגיה מאלפת: אצל מכונות לומדות, ביטחון ונכונות הם שני גדלים נפרדים, וכל אחד מהם צריך מדידה משלו.
\section{כשהדוגמאות לא מייצגות את כולם}
טעות יכולה להיות גם שיטתית. המכון האמריקאי לתקנים ולטכנולוגיה (NIST) מונה, במסגרת לניהול סיכוני AI שפרסם בינואר 2023, שלוש קטגוריות של הטיה: מערכתית, חישובית-סטטיסטית ואנושית-קוגניטיבית. לפי המסמך, כל אחת מהן יכולה להתקיים גם בלי דעה קדומה ובלי כוונה להפלות, ההטיות החישוביות נובעות לעיתים קרובות ממדגמים לא מייצגים, ומערכות AI יכולות להגדיל את המהירות ואת ההיקף של הטיות \cite{S-ch06-10}.
כפי שראינו בפרק 2, מחקר Gender Shades מ-2018 מצא ששלושה מסווגי מגדר מסחריים שגו אצל נשים כהות עור בשיעור של 20.8\% עד 34.7\%, לעומת 0.0\% עד 0.8\% אצל גברים בהירי עור \cite{S-ch02-01}. מה שגרם לפער לא הוכח: החוקרות רק משערות שהוא קשור לתת-ייצוג בנתוני האימון, שהרי החברות לא פרסמו על מה אומנו המסווגים \cite{S-ch02-01}. וגם התחום אחר: מסווגי תמונות, לא מודלי שפה. לכן הקישור ל"מדגמים לא מייצגים" של NIST הוא פרשנות של הספר.
\section{ומה קרה למה שהיא הקלידה?}
ובחזרה לשאלה השנייה של הקוראת. מודל לומד מדוגמאות, ולפעמים הוא לומד אותן בעל-פה. קרליני ועמיתיו הראו שאפשר לשחזר דוגמאות אימון בודדות מתוך מודל שפה רק באמצעות שאילתות אליו, והדגימו זאת על GPT-2 \cite{S-ch06-06}. בין הרצפים שחולצו מילה במילה היו שמות, מספרי טלפון וכתובות דוא"ל, וגם קוד, וההתקפה הצליחה אפילו ברצפים שהופיעו במסמך אחד בלבד בנתוני האימון \cite{S-ch06-08}. באחת הדוגמאות, בעקבות קידומת טקסט, השלים המודל שם, דוא"ל, טלפון, פקס וכתובת של אדם אמיתי, והחוקרים השחירו את הפרטים כדי להגן עליו \cite{S-ch06-08}. לפי התקציר, מודלים גדולים פגיעים יותר מקטנים \cite{S-ch06-06}.
\begin{behindbox}{איך תופסים מודל שמשנן}
בגרסת הכנס של המחקר (2021) הפיקו החוקרים מ-GPT-2 טקסטים רבים, ובדקו אם הם מועתקים מילה במילה מנתוני האימון \cite{S-ch06-08}:
\begin{enumerate}
\item מתוך הטקסטים נבחרו 1,800 מועמדים.
\item כל מועמד נבדק מול נתוני האימון. נמצאו 604 דוגמאות משוננות ייחודיות, כלומר 33.5\%; בווריאנט המוצלח ביותר של השיטה, 67\% \cite{S-ch06-08}.
\item הדוגמאות סווגו ביד. בין השאר, 46 כללו שמות של אנשים אמיתיים (מחוץ לטקסטים חדשותיים) ו-32 כללו פרטי קשר \cite{S-ch06-08}.
\end{enumerate}
המחברים עצמם ממתנים את התמונה: המודל GPT-2 אומן על נתונים ציבוריים, ולכן ההתקפה אינה חמורה במיוחד; על מודל שאומן על מאגר רגיש, לדבריהם, התוצאות היו חזקות יותר \cite{S-ch06-08}.
\end{behindbox}
המסגרת של NIST מזהירה שמערכות AI עלולות ליצור סיכוני פרטיות חדשים, באמצעות הסקה שמזהה אנשים או חושפת עליהם מידע פרטי, ומציעה כלים כמו טכנולוגיות לשיפור פרטיות, הסרת פרטים מזהים וצבירה של נתונים \cite{S-ch06-10}. האם השאלה של הקוראת תשמש לאימון מודל עתידי? זה תלוי בשירות שבו היא משתמשת ובתנאים שלו, ובדיוק את זה כדאי לברר לפני שמקלידים.
\begin{thinkbox}
באותה מסגרת מתוארת פרטיות כנורמות ומעשים שמגינים על האוטונומיה, הזהות והכבוד של האדם \cite{S-ch06-10}.
דמיינו שאתם מקלידים למערכת שאלה שכוללת שם של חברה, פרט רפואי או מסמך מהעבודה. על האוטונומיה, הזהות או הכבוד של מי השפיעה ההקלדה? האם הייתם מוסרים את אותו משפט לאדם זר? ומי אחראי לגורלו: מי שהקליד, מי שמפעיל את השירות, או שניהם?
\end{thinkbox}
\section{של מי הטקסט?}
גם כאן יש למכון מה לומר: נתוני אימון עשויים להיות כפופים לזכויות יוצרים, ויש לנהוג בהם לפי דיני הקניין הרוחני החלים \cite{S-ch06-10}. משרד זכויות היוצרים של ארצות הברית מפרסם דוח בשלושה חלקים: על העתקים דיגיטליים (31 ביולי 2024), על זכאות להגנה (29 בינואר 2025), ועל אימון של מערכות AI יוצרות, בגרסה טרום-סופית מ-9 במאי 2025 \cite{S-ch06-11}. לחלק השני קדמו יותר מ-10,000 הערות ציבור \cite{S-ch06-12}. המסקנה שלו: הדוקטרינות הקיימות מספיקות, וההכרעה נעשית מקרה למקרה. תוצר יכול להיות מוגן, כולו או בחלקו, כשה-AI משמש ככלי ואדם קבע את יסודות ההבעה, אבל הנחיות טקסט (prompts) לבדן, בשלב זה, כנראה אינן מספיקות; והמשרד אינו ממליץ, בשלב זה, על שינוי חקיקה \cite{S-ch06-12}.
כל זה דין אמריקאי, ואין ללמוד ממנו על המצב בישראל או באירופה. ואת מסקנות החלק השלישי, על האימון עצמו, הספר לא בחן.
\section{שימוש אחראי}
לפי NIST, למערכת AI אמינה יש שבעה מאפיינים, בהם בטיחות, אחריותיות ושקיפות, הגנה על פרטיות והוגנות, והתוקף והאמינות משמשים בסיס לכל השאר \cite{S-ch06-10}. המסמך מודה שיש ביניהם שיקולי איזון: טכניקות שמשפרות פרטיות, למשל, עלולות לפגוע בדיוק כשהנתונים דלילים \cite{S-ch06-10}. ובמקום שבו המערכת אינה יכולה לזהות או לתקן את שגיאותיה, ייתכן שיידרש אדם \cite{S-ch06-10}.
הקוראת חוזרת אל המשפט ההוא. היא לא מוחקת את התשובה כולה, אבל גם לא מקבלת אותה כמות שהיא. ביטחון עצמי אינו ראיה, והראיה, אם יש כזו, נמצאת מחוץ לחלון הצ'אט. בפרק האחרון נראה מה קורה כשהיא מתחילה לבדוק: האדם שבמעגל.
\begin{figure}[htbp]
\centering
% Figure 6.1: three case cards. Every number appears in chapter 6 / chapter 2 and in the cited source.
\newcommand{\casecard}[6]{% #1 x, #2 heading, #3 what happened, #4 what the source says, #5 lesson, #6 source
  \node[draw=midnight,line width=1pt,rounded corners=2mm,fill=white,minimum width=5.1cm,minimum height=8.2cm,anchor=north] at (#1,0) {};
  \node[text width=4.6cm,align=right,font=\headingfont\small,text=turquoise,fill=midnight,rounded corners=1.5mm,
        inner sep=2mm,anchor=north] at (#1,-0.15) {#2};
  \node[text width=4.6cm,align=right,font=\footnotesize,anchor=north] at (#1,-1.15)
       {{\headingfont מה קרה:} #3\par\medskip{\headingfont מה המקור אומר:} #4\par\medskip
        {\color{violet}\headingfont מה לומדים:} #5};
  \draw[midnight!30] (#1-2.4,-7.35) -- (#1+2.4,-7.35);
  \node[font=\scriptsize,anchor=north] at (#1,-7.45) {#6};}
\newcommand{\hl}[1]{\colorbox{gold!45}{\strut #1}}
\begin{bookdiagram}[x=1cm,y=1cm]
  \casecard{11.0}{Mata v. Avianca, 2023}
    {עורכי דין הגישו לבית המשפט החלטות שיפוטיות שאינן קיימות, שנוצרו בכלי ChatGPT.}
    {בין היתר מנה בית המשפט בשמן \hl{שש} החלטות כאלה. הוטל קנס של \hl{5,000} דולר.}
    {אין פסול בעצם השימוש בכלי. עורכי הדין הם שומרי הסף, והכשל היה בוויתור על הבדיקה.}
    {\cite{S-ch06-13}, עמ' 1, 34}
  \casecard{5.6}{Carlini et al., 2021 (GPT-2)}
    {החוקרים הפיקו טקסטים מ-GPT-2 ובדקו אם הועתקו מילה במילה מנתוני האימון.}
    {מתוך \hl{1,800} מועמדים נמצאו \hl{604} דוגמאות משוננות ייחודיות.}
    {מודל יכול לשנן נתוני אימון ולחשוף אותם. סייג: GPT-2 אומן על נתונים ציבוריים.}
    {\cite{S-ch06-08}, עמ' 2640, 2641}
  \casecard{0.2}{Gender Shades, 2018 \textbullet\ ראו פרק 2}
    {שלושה מסווגי מגדר מסחריים נבדקו בנפרד לפי קבוצות.}
    {שגיאה של \hl{20.8\%–34.7\%} אצל נשים כהות עור, מול \hl{0.0\%–0.8\%} אצל גברים בהירי עור.}
    {הסיבה לפער היא השערה ולא הוכחה, ומדובר במסווגי תמונות ולא במודלי שפה.}
    {\cite{S-ch02-01}, טבלה 4, עמ' 9}
\end{bookdiagram}

\caption{שלושה תיקים על השולחן}\label{fig:6.1}
\figurenote{מקור: \cite{S-ch06-13} (עמ' 1, 34), \cite{S-ch06-08} (עמ' 2640, 2641), \cite{S-ch02-01} (טבלה 4, עמ' 9); אין בכרטיסים נתון שאינו במקורות}
\end{figure}
\begin{insightsbox}
\begin{itemize}
\item הזיה היא תוכן שסותר את המקור או שאי אפשר לאמת מתוכו; היא נולדת גם בנתונים וגם בשלבי האימון והפענוח \cite{S-ch06-02}.
\item שטף ניסוח אינו אימות. בתיק Mata v. Avianca הכלי "אישר" תקדימים בדויים כשנשאל עליהם \cite{S-ch06-13}, ולכן שאלת "האם זה נכון?" למערכת עצמה אינה בדיקה.
\item אצל רשתות סיווג מודרניות נמצא כיול גרוע, ובדוגמה המרכזית הביטחון עלה בהרבה על הדיוק \cite{S-ch06-04,S-ch06-05}; לגבי מודלי שפה זו אנלוגיה, לא ממצא.
\item הטיה אינה דורשת כוונה רעה, ומדגמים לא מייצגים הם מקור נפוץ שלה \cite{S-ch06-10}; פערי שגיאה בין קבוצות אפשר למדוד גם כשהסיבה עצמה לא הוכחה \cite{S-ch02-01}.
\item מודל עלול לשנן ולפלוט פרטים אישיים מנתוני האימון \cite{S-ch06-08}, ולכן מה שמקלידים אינו עניין פרטי בהכרח.
\item השימוש בכלי אינו פסול; האחריות לבדוק נשארת אצל האדם \cite{S-ch06-13}.
\end{itemize}
\end{insightsbox}
